%% NOTE: Submission must be anonymized
%%       ***Check before submitting***

\documentclass[manuscript,review,anonymous]{acmart}
% \documentclass[manuscript,review]{acmart}
% \documentclass[sigconf, authordraft]{acmart}



% %% Rights management information.  This information is sent to you
% %% when you complete the rights form.  These commands have SAMPLE
% %% values in them; it is your responsibility as an author to replace
% %% the commands and values with those provided to you when you
% %% complete the rights form.
% \setcopyright{acmlicensed}
% \copyrightyear{2018}
% \acmYear{2018}
% \acmDOI{XXXXXXX.XXXXXXX}
% %% These commands are for a PROCEEDINGS abstract or paper.
% \acmConference[Conference acronym 'XX]{Make sure to enter the correct
%   conference title from your rights confirmation email}{June 03--05,
%   2018}{Woodstock, NY}
% %%
% %%  Uncomment \acmBooktitle if the title of the proceedings is different
% %%  from ``Proceedings of ...''!
% %%
% %%\acmBooktitle{Woodstock '18: ACM Symposium on Neural Gaze Detection,
% %%  June 03--05, 2018, Woodstock, NY}
% \acmISBN{978-1-4503-XXXX-X/2018/06}

%% CHI 2026 Papers - submission 3265
\acmSubmissionID{3265}



\theoremstyle{definition}
\newtheorem{definition}{Definition}

\theoremstyle{plain}
\newtheorem{theorem}{Theorem}

\begin{document}

\title[Theory of Anticipation]{Formulation of Engagement as Local Anticipation}

\author{Jeeheon (Lloyd) Oh}
\orcid{0009-0007-3192-1063}
\affiliation{%
  \department{Research Division}
  \institution{Farseer Co., Ltd.}
  \city{Seoul}
  \country{Republic of Korea}}
\email{aka.louis@outlook.com}

\begin{abstract}
  Quantifying player engagement (often called ``fun'') in game design remains a significant challenge.
  The traditional Flow theory emphasizes the balance between challenge and skill but provides qualitative rather than quantitative guidance.
  As a result, it relies on subjective evaluation methods that cannot predict engagement levels or establish mathematical optimality conditions.
  We introduce a novel metric, Local Anticipation, which measures the probability-weighted standard deviation of outcome desirabilities for a given game state, capturing anticipation of meaningful future events.
  Derived from observed principles, the formulation yields a theoretical bound $A(s) \le 0.5$ for binary-payoff games, with the coin toss achieving this maximum $(A=0.5)$ and proven optimal for single-turn engagement.
  In our experiments, we extended HpGame, a simplified one-on-one fighting game, to achieve a higher Game Design Score (GDS) and identified the parameter value that yields the highest GDS.
  This study establishes a rigorous, quantitative foundation for measuring and optimizing player engagement.
\end{abstract}

\begin{CCSXML}
  <ccs2012>
    <concept>
      <concept_id>10010405.10010476.10011187.10011190</concept_id>
      <concept_desc>Applied computing~Computer games</concept_desc>
      <concept_significance>500</concept_significance>
    </concept>
    <concept>
      <concept_id>10002950.10003714.10003716</concept_id>
      <concept_desc>Mathematics of computing~Mathematical optimization</concept_desc>
      <concept_significance>300</concept_significance>
    </concept>
    <concept>
      <concept_id>10003120.10003121.10011748</concept_id>
      <concept_desc>Human-centered computing~Empirical studies in HCI</concept_desc>
      <concept_significance>100</concept_significance>
    </concept>
  </ccs2012>
\end{CCSXML}

\ccsdesc[500]{Applied computing~Computer games}
\ccsdesc[300]{Mathematics of computing~Mathematical optimization}
\ccsdesc[100]{Human-centered computing~Empirical studies in HCI}

\keywords{Game Theory, Quantitative Game Design, Player Engagement, Decision Theory, Mathematical Optimization, Entertainment Computing}

\maketitle

\section{Introduction}
Quantifying player engagement in game design remains a fundamental challenge.
Established theories such as Flow theory~\cite{csikszentmihalyi1990flow} provide qualitative frameworks for understanding optimal experience but lack the mathematical precision required for systematic design optimization, and the objective comparison of game mechanics.

This limitation constrains design innovation, forcing developers to rely on intuition due to the lack of rigorous methods to evaluate novel concepts.
Existing evaluation techniques depend on subjective assessments---which cannot establish mathematical optimality---or require full implementation for playtesting, making systematic design exploration costly and inefficient.
Consequently, the absence of objective engagement metrics forces developers to favor established patterns over innovative mechanics.

To address this challenge, we introduce \emph{local anticipation}, a mathematical framework that quantifies engagement as the probability-weighted standard deviation of outcome desirabilities---a metric computable directly from game state definitions.

We demonstrate the analytical power of this framework by providing a complete mathematical characterization of single-turn games, deriving fundamental engagement bounds, and proving that the coin toss achieves optimal design among single-turn games for maximizing player anticipation.

\section{Related Work}
Flow theory, introduced by Csikszentmihalyi~\cite{csikszentmihalyi1990flow}, identifies optimal experience as arising from the balance between challenge and skill levels.
While Flow theory remains the most influential theoretical framework for understanding engagement in game design, it provides only qualitative guidance and cannot quantify the role of uncertainty or distinguish between various sources of engagement within optimal flow states.
These limitations render Flow theory insufficient for systematic game design optimization.

GameFlow~\cite{sweetser2005gameflow} adapted flow theory into eight game-specific elements and demonstrated that structured criteria could differentiate between successful and failed games.
Subsequent frameworks like PLEX~\cite{arrasvuori2011plex} categorized 22 playful experiences, and the Player Experience Inventory~\cite{abeele2020pxi} provided dual-level measurement validated across multiple studies.
Despite their theoretical sophistication, these approaches remain fundamentally qualitative assessment tools rather than mathematical optimization frameworks.

Existing game metrics primarily assess player behavior (e.g., retention, progression) or subjective responses (e.g., surveys, reviews) rather than fundamental engagement properties.
While valuable for post-launch optimization, these approaches cannot inform design decisions during development and lack the mathematical rigor required for systematic design exploration.

The lack of quantitative frameworks for measuring engagement leaves game design without objective tools to evaluate novel mechanics or establish mathematical optimality conditions, creating a critical gap between theoretical understanding and practical design needs.

\section{Derivation}
The Theory of Anticipation arises from fundamental observations about the mathematical nature of engagement, demonstrating why engagement is better modeled by transition variance rather than static state properties.

\subsection{Static States Generate No Engagement}
Consider a highly desirable game state---such as a commanding lead or optimal positioning---followed by indefinitely pausing the game.
Despite occupying the most advantageous position, this scenario produces zero engagement.
This \emph{paused game paradox} demonstrates that engagement cannot arise from state desirability alone.

Similarly, engagement disappears when outcomes become certain: guaranteed victory or inevitable defeat eliminate engagement regardless of how favorable the situation appears.

\subsection{Engagement Requires Transition Uncertainty}
Engagement arises solely from uncertainty in meaningful state transitions.
Consider a penalty-only scenario: losing 10 points with 90\% probability versus losing 100 points with 10\% probability.
Although all outcomes are undesirable, the uncertainty about which penalty will occur generates measurable engagement.

%% Note: use sloppypar to fix 'Overfull \hbox' warning
\begin{sloppypar}
  This demonstrates that variance among potential outcomes drives engagement, even when all outcomes represent losses.
\end{sloppypar}

\subsection{Variance-Based Formulation}
These observations show that neither average desirability nor absolute outcome values determine engagement.
Instead, the mathematical spread among possible outcomes drives engagement, regardless of whether they represent gains or losses.

This motivates a variance-based approach: engagement should be measured by the probability-weighted dispersion of outcome desirabilities, requiring the weighted standard deviation formulation that follows.

\section{Local Anticipation}
We model a game as a set of states $S$ with transition probabilities $P(s \to s')$ and a desirability function $D: S \to \mathbb{R}$.
For any state $s \in S$, the expected desirability is:

\begin{equation}
  \mu(s) = \sum_{s' \in S} P(s \to s') \cdot D(s')
\end{equation}

The \emph{local anticipation} at state $s$ measures engagement as the probability-weighted standard deviation of outcome desirabilities:

\begin{equation}
  A(s) = \sqrt{\sum_{s' \in S} P(s \to s') \cdot \big(D(s') - \mu(s)\big)^2}
\end{equation}

This formulation captures anticipation as the uncertainty over meaningful future outcomes.
High anticipation occurs when transition probabilities are distributed among outcomes with substantially different desirabilities.

For single-turn games, we restrict to terminal states where $P(s \to s) = 1$ for all outcomes, reducing the analysis to the probability distribution over final desirabilities.

\section{Enabling Objective Analysis}
To eliminate subjectivity, we assign desire values as follows: winning states receive $D(\text{win}) = 1$, and all others receive $D(\text{other}) = 0$.
We term this the \emph{canonical intrinsic desire}.

Using this function standardizes evaluation across games, enabling objective comparison under identical conditions.
Although the desire function can be customized for experimental purposes, we found this function sufficient for most applications encountered.

\section{Theoretical Properties}
We establish some properties of the local anticipation metric:

\begin{theorem}[Anticipation Bound for Binary-Payoff Games]\label{thm:anticipation_bound_binary_payoff}
  For binary-payoff games, $A(s) \leq 0.5$ with equality achieved when $p = 0.5$.
\end{theorem}

\begin{proof}
  Consider a binary-payoff game with win probability $p$.
  Assume $D(\text{win}) = 1$ and $D(\text{lose}) = 0$.

  The expected desirability is $\mu(s) = p$, so:
  \[
    A(s) = \sqrt{p(1-p)^2 + (1-p)p^2} = \sqrt{p(1-p)}.
  \]

  To maximize $f(p) = p(1-p)$, we solve $f'(p) = 1-2p = 0$, yielding $p = 1/2$.
  Since $f''(p) = -2 < 0$, this is a maximum, with $f(1/2) = 1/4$.

  Therefore $A(s) \leq 1/2$, with equality if and only if $p = 1/2$.
\end{proof}

\begin{theorem}[Anticipation Bound for Canonical Single-Turn Games]\label{thm:anticipation_bound_canonical_single_turn}
  For any finite single-turn game under canonical desire assignment, $A(s) \leq 0.5$, with equality if and only if the probability of winning equals $0.5$.
\end{theorem}

\begin{proof}
  Under canonical desire assignment, all winning states have $D = 1$ and all non-winning states have $D = 0$, reducing the game to a binary-payoff structure.
  Consequently, the result follows directly from Theorem~\ref{thm:anticipation_bound_binary_payoff}, irrespective of the number of terminal outcomes.
\end{proof}

\section{Optimal Single-Turn Game Design}
From the theoretical bound $A \le 0.5$, we identify the fair coin toss as the optimal single-turn game for maximizing player anticipation in canonical settings.

\begin{theorem}[Optimality of the Fair Coin Toss]
  Among all canonical single-turn games, the fair coin toss uniquely achieves the maximum local anticipation $A=0.5$ \emph{if and only if} the probability of winning is $0.5$.
\end{theorem}

\begin{proof}
  By Theorem~\ref{thm:anticipation_bound_canonical_single_turn}, the maximum anticipation $A = 0.5$ is attained exactly when $p = 0.5$.
  The fair coin toss satisfies this condition with $P(\text{heads}) = P(\text{tails}) = 0.5$, and thus achieves the maximum.
  Conversely, if any canonical single-turn game achieves $A = 0.5$, then its win probability must be $p = 0.5$, which is precisely the fair coin toss case.
\end{proof}

Formally, this optimality condition corresponds to solving:
\[
  \max_{p \in [0,1]} \sqrt{p(1-p)}
\]

\section{Expansion of the Framework}
\subsection{Multi-Turn Games}
Unlike single-turn games, real games involve sequences of decisions extending beyond a single state transition.
Our framework naturally extends to multi-turn games.
For acyclic multi-turn games, expected desirability values can be computed efficiently through dynamic programming by backpropagating from terminal states:

\begin{definition}[Multi-Turn Desire Propagation]
  Given terminal states $T \subseteq S$ with known desirability values $D(s)$ for all $s \in T$, we propagate to non-terminal states in reverse topological order:
  \begin{equation}
    D(s) = \sum_{s' \in S} P(s \to s') \cdot D(s')
  \end{equation}
\end{definition}

Then we can compute the anticipation for each state using the Local Anticipation formula.
Since each state's expected desirability is calculated exactly once and each calculation involves only a fixed number of transitions, the algorithm runs in $O(|S|)$ time and scales efficiently as long as the state space is not excessively large.

\subsection{Nested Decomposition of Anticipation}
The most significant extension involves \emph{nested anticipation analysis}, where anticipation values themselves become desire functions for higher-order calculations.
This recursive structure produces Taylor-series-like anticipation components $A_1, A_2, A_3, \ldots$, each capturing different engagement depths.

\begin{definition}[Hierarchical Anticipation]
  Define the local anticipation function $\mathcal{A}(s; D)$ for state $s$ with desire function $D$:
  \begin{equation}
    \mathcal{A}(s; D) = \sqrt{\sum_{s' \in S} P(s \to s') \cdot (D(s') - \mu_D(s))^2}
  \end{equation}
  where
  \[
    \mu_D(s) = \sum_{s' \in S} P(s \to s') \cdot D(s')
  \]
  The hierarchical anticipation components are defined recursively as:
  \begin{align}
    A_1(s) &= \mathcal{A}(s; D), \\
    A_2(s) &= \mathcal{A}(s; A_1), \\
    A_n(s) &= \mathcal{A}(s; A_{n-1}) \quad \text{for } n \geq 3.
  \end{align}
\end{definition}

This decomposition captures respective depths of engagement:
\begin{itemize}
  \item $A_1$: Immediate excitement from direct outcomes,
  \item $A_2$: Longer-term anticipation from varied future possibilities,
  \item $A_{3+}$: Longer-term and deeper anticipation from complex strategic interactions.
\end{itemize}

\section{Experiments}
We applied our framework to measure the \emph{Game Design Score} (GDS) in a simple game, HpGame.
By incorporating a rage and critical hit system into HpGame, we improved the GDS and identified the optimal critical hit chance for game design.

The Game Design Score is calculated by averaging the combined anticipation components $A_1$ through $A_5$ across all game states, weighted by the reachability of each state from the initial state.

\subsection{HpGame}
HpGame models a simplified one-on-one fighting game where each player has health points and can attack with a probability of hitting or missing.
Players alternate turns until one player's health reaches zero.

\paragraph{Model Definition.}
%% Note: use allowbreak to fix 'Overfull \hbox' warning
Two players each have HP in the set $\{1, 2, 3,\allowbreak 4, 5\}$.
Each turn results in one of three equally likely outcomes:
\begin{align}
  P(\text{Player 1 hits, Player 2 misses}) &= \frac{1}{3} \\
  P(\text{Both players hit}) &= \frac{1}{3} \\
  P(\text{Player 1 misses, Player 2 hits}) &= \frac{1}{3}
\end{align}
Each successful hit deals 1 point of damage.

\paragraph{Results.}
The Game Design Score (GDS) of the HpGame model is calculated from five metrics: $A_1 = 0.176$, $A_2 = 0.135$, $A_3 = 0.072$, $A_4 = 0.030$, and $A_5 = 0.017$, summing to 0.430.

The most engaging moments occur when HP values are close and win probabilities are moderate.
The highest individual $A_1$ value occurs at the state $(1, 1)$, as shown in Table~\ref{tab:most_engaging_states}, which corresponds to a coin-toss scenario.
Note that $A_1$ does not reach the theoretical maximum of 0.5 because the win probability is not exactly 0.5 due to the possibility of mutual hits (resulting in a draw).

Interestingly, the sum of $A_1$ through $A_5$ can exceed 0.5, as observed in the top-ranked state.
This suggests potential for an unbounded conjecture regarding hierarchical anticipation components.

\begin{table}
  \caption{Top 15 most engaging game states in HpGame, ranked by sum of $A_1$--$A_5$. States are denoted as tuples (Player1 HP, Player2 HP).}
  \label{tab:most_engaging_states}
  \setlength{\tabcolsep}{4pt}
  \begin{tabular}{cccccccc}
    \toprule
    Player1 HP & Player2 HP & $A_1$ & $A_2$ & $A_3$ & $A_4$ & $A_5$ & Sum \\
    \midrule
    2 & 1 & 0.31 & 0.22 & 0.00 & 0.00 & 0.00 & 0.54 \\
    3 & 5 & 0.09 & 0.23 & 0.12 & 0.06 & 0.03 & 0.53 \\
    4 & 2 & 0.12 & 0.24 & 0.08 & 0.03 & 0.02 & 0.50 \\
    5 & 3 & 0.12 & 0.24 & 0.06 & 0.07 & 0.02 & 0.50 \\
    4 & 5 & 0.14 & 0.20 & 0.04 & 0.08 & 0.02 & 0.48 \\
    3 & 2 & 0.22 & 0.18 & 0.03 & 0.05 & 0.00 & 0.48 \\
    1 & 1 & 0.47 & 0.00 & 0.00 & 0.00 & 0.00 & 0.47 \\
    5 & 2 & 0.05 & 0.19 & 0.18 & 0.02 & 0.02 & 0.47 \\
    2 & 4 & 0.08 & 0.21 & 0.13 & 0.03 & 0.02 & 0.46 \\
    3 & 4 & 0.15 & 0.22 & 0.02 & 0.07 & 0.01 & 0.46 \\
    2 & 2 & 0.28 & 0.07 & 0.10 & 0.00 & 0.00 & 0.46 \\
    4 & 3 & 0.18 & 0.16 & 0.06 & 0.05 & 0.01 & 0.45 \\
    5 & 4 & 0.16 & 0.14 & 0.08 & 0.04 & 0.02 & 0.44 \\
    2 & 3 & 0.16 & 0.23 & 0.00 & 0.05 & 0.00 & 0.44 \\
    3 & 1 & 0.10 & 0.22 & 0.10 & 0.00 & 0.00 & 0.43 \\
    \bottomrule
  \end{tabular}
\end{table}

\subsection{HpGame\_Rage}
We extend HpGame with the novel ``rage/crit'' mechanics.
This rage/crit system yielded the highest GDS in our trials.
The system is defined as follows:
\begin{itemize}
  \item Rage accumulates when dealing or receiving damage.
  \item Critical hits deal additional damage calculated as $\text{damage} = 1 + (\text{rage} \times \text{rage\_multiplier})$, where the default multiplier is 1.
  \item Rage persists and does not get consumed by critical hits; it only increases over time.
  \item Players can land critical hits with a fixed probability (10\% in this experiment).
\end{itemize}

\paragraph{Model Definition.}
The rage system expands the transition space to 8 possible outcomes per turn, where $p_h$ denotes hit probability, $p_c$ critical hit probability, and $p_m$ miss probability.
Note that the mutual miss scenario is excluded to prevent cyclic states.
Its probability is redistributed evenly among the remaining outcomes.
\begin{align}
  &P(\text{P1 hits, P2 misses}) = p_h \cdot p_m \\
  &P(\text{P1 crits, P2 misses}) = p_c \cdot p_m \\
  &P(\text{P1 misses, P2 hits}) = p_m \cdot p_h \\
  &P(\text{P1 misses, P2 crits}) = p_m \cdot p_c \\
  &P(\text{Both hit}) = p_h \cdot p_h \\
  &P(\text{P1 hits, P2 crits}) = p_h \cdot p_c \\
  &P(\text{P1 crits, P2 hits}) = p_c \cdot p_h \\
  &P(\text{Both crit}) = p_c \cdot p_c
\end{align}

\paragraph{Results.}
TODO:
The Game Design Score of the HpGame model is calculated from five metrics: $A_1 = 0.xxx$, $A_2 = 0.xxx$, $A_3 = 0.xxx$, $A_4 = 0.xxx$, and $A_5 = 0.xxx$, summing to 0.xxx.

We observe a significant increase in GDS, improving from 0.430 to 0.544---a 26.5\% improvement.
The rage system notably boosts higher-order components like $A_4$ and $A_5$, reflecting enhanced multi-turn anticipation depth.

The state $(4, 2, 3, 3)$ achieves a total anticipation of 0.94 (sum of components: 0.26 + 0.31 + 0.22 + 0.11 + 0.04), as shown in Table~\ref{tab:most_engaging_states_rage}.
This state corresponds to an effective match point where both players can potentially eliminate each other with one critical hit.
This creates immediate tension alongside multiple branching possibilities, producing a balanced mix of action-packed immediacy and long-term strategic depth.

Similarly, the state $(5, 2, 3, 2)$ is notable.
Although the player with 5 HP faces lower immediate risk (hence slightly lower $A_1$), a single critical hit can dramatically shift the game balance.
This state illustrates potential unpredictable turn-overs and branching outcomes over the long term, which is reflected in higher $A_4$ and $A_5$ values.

It is noteworthy that a slight asymmetry in balance tends to create generally more engaging states.
When a player is slightly ahead or behind and a single event can swing the balance, each action feels more consequential---winning the exchange means gaining the upper hand, fostering a sense of momentum.

These observations provide strong empirical support for our theoretical framework.
As shown in the hierarchical anticipation decomposition, higher-order components effectively capture narrative elements such as dramatic swings, strategic branches, and long-term tension arcs.
The rage system's ability to naturally elevate these components demonstrates how well-designed mechanics can enhance engagement across multiple temporal scales simultaneously.

\begin{table}
  \caption{Top 15 most engaging game states in HpGame with rage mechanics, ranked by sum of $A_1$--$A_5$. States are denoted as tuples (Player1 HP, Player1 Rage, Player2 HP, Player2 Rage).}
  \label{tab:most_engaging_states_rage}
  \setlength{\tabcolsep}{4pt}
  \begin{tabular}{cccccccccc}
    \toprule
    \multicolumn{2}{c}{Player1} & \multicolumn{2}{c}{Player2} & \multicolumn{6}{c}{} \\ 
    \cmidrule(lr){1-2} \cmidrule(lr){3-4}
    HP & Rage & HP & Rage & $A_1$ & $A_2$ & $A_3$ & $A_4$ & $A_5$ & Sum \\
    \midrule
    4 & 2 & 3 & 3 & 0.26 & 0.31 & 0.22 & 0.11 & 0.04 & 0.94 \\
    4 & 3 & 3 & 3 & 0.26 & 0.31 & 0.22 & 0.11 & 0.04 & 0.94 \\
    3 & 3 & 4 & 3 & 0.27 & 0.30 & 0.21 & 0.11 & 0.03 & 0.91 \\
    3 & 3 & 4 & 2 & 0.27 & 0.30 & 0.21 & 0.11 & 0.03 & 0.91 \\
    3 & 4 & 3 & 4 & 0.27 & 0.29 & 0.17 & 0.06 & 0.02 & 0.82 \\
    3 & 3 & 3 & 4 & 0.27 & 0.29 & 0.17 & 0.06 & 0.02 & 0.82 \\
    3 & 3 & 3 & 2 & 0.27 & 0.29 & 0.17 & 0.06 & 0.02 & 0.82 \\
    3 & 3 & 3 & 3 & 0.27 & 0.29 & 0.17 & 0.06 & 0.02 & 0.82 \\
    3 & 2 & 3 & 3 & 0.27 & 0.29 & 0.17 & 0.06 & 0.02 & 0.82 \\
    3 & 4 & 3 & 3 & 0.27 & 0.29 & 0.17 & 0.06 & 0.02 & 0.82 \\
    4 & 2 & 3 & 2 & 0.23 & 0.26 & 0.19 & 0.10 & 0.04 & 0.81 \\
    4 & 3 & 3 & 2 & 0.23 & 0.26 & 0.19 & 0.10 & 0.04 & 0.81 \\
    5 & 2 & 3 & 2 & 0.18 & 0.24 & 0.21 & 0.13 & 0.06 & 0.81 \\
    4 & 2 & 2 & 3 & 0.28 & 0.25 & 0.16 & 0.06 & 0.02 & 0.77 \\
    4 & 4 & 2 & 4 & 0.28 & 0.25 & 0.16 & 0.06 & 0.02 & 0.77 \\
    \bottomrule
  \end{tabular}
\end{table}

\subsection{Parameter Optimization}
We optimized the critical hit chance in HpGame\_Rage to maximize engagement.
Optimization was performed by a simple brute-force search over 1\% increments from 0\% to 40\%.

\paragraph{Results.}
We found that the optimal critical hit chance was 13\%, exhibiting a clear inverted-U relationship across the parameter space, as shown in Table~\ref{tab:critical_hit_optimization}.
This pattern is reasonable:
when the critical hit chance is too low, the game behavior converges toward the baseline HpGame without additional mechanics;
conversely, when the critical hit chance is too high, small tactical advantages (e.g., 1-HP differences) become less relevant, as critical hits increasingly dominate game outcomes.

\begin{table}
  \caption{Critical hit probability optimization results. (Bold: optimal critical hit chance)}
  \label{tab:critical_hit_optimization}
  \begin{minipage}[t]{0.45\columnwidth}
    \raggedleft
    \begin{tabular}{cc}
      \toprule
      Critical \% & GDS \\
      \midrule
      0\% & 0.430 \\
      1\% & 0.434 \\
      2\% & 0.451 \\
      3\% & 0.469 \\
      4\% & 0.485 \\
      5\% & 0.500 \\
      6\% & 0.513 \\
      7\% & 0.523 \\
      8\% & 0.532 \\
      9\% & 0.539 \\
      10\% & 0.544 \\
      11\% & 0.548 \\
      12\% & 0.550 \\
      \textbf{13\%} & \textbf{0.551} \\
      14\% & 0.551 \\
      15\% & 0.549 \\
      16\% & 0.548 \\
      17\% & 0.545 \\
      18\% & 0.542 \\
      19\% & 0.538 \\
      20\% & 0.534 \\
      \bottomrule
    \end{tabular}
  \end{minipage}
  \hfill
  \begin{minipage}[t]{0.45\columnwidth}
    \raggedright
    \begin{tabular}{cc}
      \toprule
      Critical \% & GDS \\
      \midrule
      21\% & 0.529 \\
      22\% & 0.524 \\
      23\% & 0.519 \\
      24\% & 0.514 \\
      25\% & 0.509 \\
      26\% & 0.503 \\
      27\% & 0.498 \\
      28\% & 0.493 \\
      29\% & 0.487 \\
      30\% & 0.482 \\
      31\% & 0.476 \\
      32\% & 0.471 \\
      33\% & 0.466 \\
      34\% & 0.461 \\
      35\% & 0.456 \\
      36\% & 0.451 \\
      37\% & 0.446 \\
      38\% & 0.441 \\
      39\% & 0.436 \\
      40\% & 0.432 \\
      & \\
      \bottomrule
    \end{tabular}
  \end{minipage}
\end{table}

This case provides practical validation of computer-assisted parameter optimization methodology in game design.
While intuition might reasonably suggest that 10--15\% is a viable critical hit range based on conventional design experience, precisely identifying 13\% as the optimal value clearly surpasses intuition-based exploration.

\section{Discussion and Future Work}
\subsection{The Game Ending Button Problem}
Consider a thought experiment: players are given a button they can press at any moment to instantly end the game with a 50\% chance of victory.
This design achieves the theoretical maximum $A_1 = 0.5$ across all game states, as every state effectively becomes a coin-toss game.

However, despite maximizing first-order anticipation, such a game would be profoundly unengaging.
This paradox illustrates why optimizing $A_1$ alone is insufficient for creating compelling game experiences.

The Game Ending Button fails because it eliminates all higher-order anticipation components.
More intuitively, it lacks strategic depth, as only the button press matters, effectively reducing the game to a series of independent single-turn coin tosses.
In short, the game lacks meaningful ``depth.''

The higher-order components effectively capture this notion of ``depth.''
Preliminary experiments indicate that the sum $(A_1 + A_2 + A_3 + \ldots)$ correlates more closely with the perceived quality of a game design than $A_1$ alone.
This is reasonable because human intuition naturally and automatically builds predictive models to some extent.
Therefore, using the combined sum of anticipation components provides a better measure of perceived anticipation than relying solely on $A_1$.

\subsection{More Standardization for the Game Design Score}
The Game Ending Button Problem highlights that multi-component analysis is essential for a meaningful measurement of engagement.
Internally, we used the sum of $A_1$ through $A_5$; however, there is currently no established methodology for optimally weighting these components or for determining the appropriate number of components necessary for standardized comparison across different games.

Are weightings necessary for the components?
How many components should be considered sufficient for a reliable standardized comparison?

Further research is required to address these questions.

\subsection{Conjecture: Higher-Order Components May Capture Replayability}
Consider a common phenomenon: a game may appear very fun initially with a high anticipation $A_1$, yet gradually become boring over time despite unchanged mechanics.
This phenomenon is often taken as evidence that ``fun'' is inherently subjective and therefore difficult to quantify.

However, the hierarchical structure of anticipation offers an elegant explanation.
Engagement decay over time occurs because, as players’ understanding of the game deepens, they become better at predicting outcomes, reducing the perceived variance of significant in-game events within a session.

The component $A_2$ can be interpreted as anticipation of anticipation, capturing perceived variance after a player has mastered $A_1$-level patterns.

From this perspective, we can predict several properties.
As players successfully model the game and master the $A_1$ hierarchy, $A_2$ becomes the primary driver of engagement.
When $A_2$ patterns become predictable, $A_3$ takes precedence, and so forth.
Games with strong higher-order components are able to maintain engagement longer through multiple mastery phases, which helps explain why strategically complex games like \textit{League of Legends} exhibit exceptional longevity compared to action-focused games with shallower strategic depth.

While this interpretation aligns well with empirical observations, further systematic analysis is needed to establish the precise relationship between higher-order anticipation components and replayability.
For example, is this relationship logarithmic or linear? What is the typical time needed for players to master each level of the anticipation hierarchy?

\subsection{Conjecture: Optimizing Higher-Order Components May Generate Narrative Structure}
Optimizing higher-order anticipation components naturally creates trade-offs with immediate engagement.
Maximizing $A_2$ requires variance in $A_1$ values.
Since $A_1$ has an upper bound of 0.5, generating non-zero variance mathematically necessitates that some $A_1$ values fall below 0.5.

This constraint suggests designing games with alternating high- and low-anticipation states, producing rhythmic engagement patterns.
Interestingly, this mathematical requirement appears to align with established storytelling principles.

Although not yet rigorously proven, our informal experiments suggest that optimization of multidimensional anticipation may converge on classical narrative structures, such as Freytag's Pyramid or the Three-Act Structure.
Further investigation of this conjecture could yield valuable insights.

\subsection{Conjecture: Anticipation May Be Unbounded for Higher-Order Components}
While $A_1 \leq 0.5$ for binary games, higher-order components may be unbounded.
This conjecture arises from the observation that deeper engagement often requires temporal development.
Creating profound emotional impact within short time frames appears fundamentally more difficult than when longer durations are available.

Following this pattern, although $A_1$ remains constrained to 0.5, the sum $A_1 + A_2$ may have a higher bound, $A_1 + A_2 + A_3$ may exceed $A_1 + A_2$, and so on.
If this conjecture holds true, we could theoretically build games with arbitrarily high anticipation by extending game length and depth, limited only by session time and event density.

Validating this conjecture would have significant implications for the game design industry.
It would suggest that, contrary to the current impression that game designs are saturated and innovation is nearly exhausted, it is possible to create games with arbitrarily high anticipation through newer methodologies such as computer-assisted design discovery.

\subsection{Computer-Assisted Design Discovery}
Components beyond $A_2$ lie beyond the domain where human intuition can effectively optimize.
Optimizing $A_2$ necessarily requires sacrificing $A_1$ values, resulting in a complex, combinatorial, multi-dimensional optimization problem that exceeds human cognitive capacity to simultaneously consider multiple competing objectives.

This aligns with the observed phenomenon that designing games with strategic depth and long-term engagement is significantly more difficult than creating immediately exciting and action-packed experiences.

Computer-assisted approaches offer systematic solutions to these limitations.
By defining formal game models and employing techniques such as genetic algorithms, gradient descent, and other optimization methods, we can discover novel game designs and identify optimal parameter configurations that navigate these complex trade-off spaces more effectively than human intuition alone.

As demonstrated in the previous sections, dynamic programming algorithms can compute nested components in $O(|S|)$ time per component.
Therefore, optimization of game design using these methods is both practical and feasible.

\section{Conclusion}
We have demonstrated that HpGame\_Rage outperforms HpGame by 26.5\%, supported by a mathematically validated comparison (GDS 0.544 vs. 0.430).
To our knowledge, this may be the first instance where two game designs have been compared with objective mathematical evidence at such a high level of certainty.

Using our framework, we successfully developed and validated novel ``rage/crit'' mechanics for 1v1 fighting games with unprecedented design certainty---eliminating the need for extensive playtesting or full implementation.

Regarding nested anticipation, we observe that higher-order components ($A_4$, $A_5$) in HpGame\_Rage effectively capture long-term strategic branches and thus narrative structure, providing empirical support for our nested anticipation decomposition approach.

Additionally, our state space analysis revealed new insights, such as the finding that ``slightly asymmetric game states generate superior engagement compared to perfectly balanced scenarios.''

Finally, we optimized the critical hit chance parameter, improving the Game Design Score from 0.544 to 0.551 (+1.3\%) by adjusting the value from 10\% to 13\%.
This serves as a practical demonstration of applying our framework to optimize game mechanics and parameters, yielding directly actionable design insights.

Achieving comparable optimization without this framework would require extensive large-scale statistical studies involving hundreds of players and fully playable prototypes---clearly impractical for the multitude of micro-design decisions in modern game development.

This positions our framework as a potentially transformative tool for game design, providing a level of mathematical and objective design insight previously unattainable with such precision.

In today's saturated gaming market, where innovation often relies on incremental genre mixing, our framework offers substantial potential to overcome current design limitations by providing systematic methods to discover and validate truly novel mechanics.

\section{Limitations and Scope}
This framework is best suited for games with well-defined objectives.  
It works well for competitive games with clear win conditions, especially competitive MOBAs, but applies poorly to games like Minecraft, where exact win conditions do not exist.

Additionally, real games operate in real-time with prohibitively large state spaces, requiring idealized and simplified models for computational feasibility.  
A key challenge is the ambiguity in defining a universally agreeable desire input, $D$, across different genres; this was straightforward for games with trivial win conditions.

Another common criticism the framework must address is subjectivity.  
For instance, a player's personal preferences, which may seem arbitrary, can significantly alter their engagement with a game.

Despite these challenges, the theory provides a robust framework with significant practical applications and deeper implications.
The primary practical use of the Theory of Anticipation (ToA) has been in inventing novel game mechanics, making objective design choices, and gaining insights beyond common intuition, rather than serving as a tool for top-down analysis of entire games.

This is comparable to the relationship between classical physics and the real world.  
While it is not practical to apply classical physics to real-world systems at full resolution, its principles guide engineering and design and work effectively in idealized models.  
Similarly, ToA enables a more scientific approach to game design.

The issue of subjectivity can be modeled as per-player random noise on the desire function, $D$.  
From this perspective, ToA provides the theoretical, averaged target for engagement.  
Therefore, counterintuitively, the existence of subjectivity does not weaken the theory's role; it highlights its value.  
The framework offers a noise-free theoretical target metric that was previously impossible to measure without large-scale, often prohibitively expensive methods such as extensive playtesting.

At present, the theory shows promising initial signs as a unifying framework.  
At its core, ToA is a single formula which, when expanded into a complex system, produces emergent properties that connect different fields, suggesting that concepts like narrative fun and gaming fun may stem from the same underlying principles.

We have also actively attempted to challenge the theory, especially by searching for cases of engagement it cannot explain.  
For example, we could not find a case where engagement arises without meaningful branching outcomes, as demonstrated in the ``Paused Game Problem'' thought experiment.  
To date, all issues encountered have been related to defining the inputs, not fundamental flaws in the theory.  
Nonetheless, further scrutiny is required to fully validate its robustness.

\section{Supplementary Materials}
The source code and browser-playable demos for the HpGame and HpGame\_Rage experiments are available in the online version of this paper.
\anon[
  Additional information has been omitted to preserve anonymity.
]{%
  Additionally, the ToA project page, which includes the experiments presented in this paper, is available at \url{https://akalouis.github.io/}.
}

\begin{acks}
  The author would like to thank Taeyeong Kim for valuable assistance in preparing the manuscript.
\end{acks}

\bibliographystyle{ACM-Reference-Format}
\bibliography{references}

\end{document}
\endinput
